\section{Exact order-constrained depolarization in arbitrary dimension}\label{sec:order}
Let \(\E_A(X)=I_d/d\otimes\Tr_AX\) and
\(\N_\lambda=\lambda\id+(1-\lambda)\E_A\), where
\(-1/(d^2-1)\le\lambda\le1\), the known physical range~\cite[Section~8.1]{HRF}. For \(\kappa>1\) restrict the finite pairs to
\begin{equation}\label{eq:order-class}
 \rho_R=\sigma_R,\quad \rho,\sigma>0,\quad \rho\ne\sigma,\quad
 \rho\preceq\kappa\E_A(\rho),\quad \sigma\preceq\kappa\E_A(\sigma).
\end{equation}
For the metric impose the base order bound and \(\Tr_AD=0\); one may also require a two-sided small neighborhood satisfying \eqref{eq:order-class}. The suprema proved below are unchanged. The budget uses the \emph{actual} reference marginal, not an optimized conditional min-entropy reference.

\begin{proposition}[Commuting Bell lower and its range]\label{prop:bell}
For every \(d\ge2\), \(1<\kappa\le d^2/2\) and every nonaffine operator-convex Petz generator normalized by \(f(1)=0\), the complete pair coefficient under \eqref{eq:order-class} is
\(\lambda^2\kappa/[1+\lambda(\kappa-1)]\) throughout the physical depolarizing range. The SLD and minimal-chi-square coefficients have the same value. A reference of dimension \(d\) suffices for the lower supremum.
\end{proposition}
\begin{proof}
The order upper is proved below. Write \(n=d^2\), \(B=I_n/n\), and take two orthogonal maximally entangled Bell projectors \(P_0,P_1\). For \(1<K<\kappa\le n/2\), set
\[
 a=K/n,\quad b=(1-2a)/(n-2),\quad
 \sigma_K=bI_n+(a-b)(P_0+P_1),\quad D=P_0-P_1.
\]
Here \(a>b>0\), \(\E_A(\sigma_K)=B\), \(\Tr_AD=0\), and the actual order constant is \(K\). The input and output commute; their eigenvalues on the two changed rays are \(a\) and \(a_\lambda=\lambda a+(1-\lambda)/n>0\). Hence the SLD ratio is exactly \(\lambda^2a/a_\lambda\).

For \(0<|t|<\min(a,b,(\kappa-K)/n)\), \(\rho_t=\sigma_K+tD\) and \(\sigma_K\) are faithful with identical reference and both order bounds. Their Petz divergences are
\[
 a\left[f(1+t/a)+f(1-t/a)\right],\qquad
 a_\lambda\left[f(1+\lambda t/a_\lambda)+f(1-\lambda t/a_\lambda)\right].
\]
For a normalized nonaffine operator-convex generator, \(f''(1)>0\), as follows from the nonzero positive representation in \eqref{eq:representation}. The Taylor ratio as \(t\to0\) is therefore \(\lambda^2a/a_\lambda=\lambda^2K/[1+\lambda(K-1)]\). Then let \(K\uparrow\kappa\). At the endpoint \(\kappa=n/2\), each earlier base remains faithful. The cases \(\lambda=0,1\) are immediate. The quadratic chi-square equality follows directly from the metric ratio.
\end{proof}

This is a reparametrization of the Bell family already present in our earlier component work, not an independent invention of Bell-assisted contraction. In particular, when \(d\ge3\), the interval \(\kappa\le d+1\) is contained in \(\kappa\le d^2/2\); the rank-\(d\) ladder previously considered on that interval is an alternative proof. A conditional-expectation index alone does not imply sharpness: the pinching \(\E=(\id+\Ad_Z)/2\) has index two, but its equal-\(\E\) contraction is \(\lambda^2\), rather than the larger raw order comparison, for \(0<\lambda<1\).

\begin{theorem}[Principal-rich rank-two order profile]\label{thm:order}
For \(d\ge2\), set \(M_d=d(2d-1)/2\). For \(1<\kappa\le M_d\), the complete SLD, minimal-chi-square and finite triangular-Petz contraction coefficients under \eqref{eq:order-class} equal
\begin{equation}\label{eq:profile}
 C_\kappa(\lambda)=\frac{\lambda^2\kappa}{1+\lambda(\kappa-1)}.
\end{equation}
The complete supremum is over finite reference systems; a reference of dimension \(2d-1\) supplies the lower. The same expression is an upper for each normalized nonaffine operator-convex Petz generator, sharp uniformly over that class. Equality for each generator outside the Bell range, or for every fixed reference dimension, is not asserted.
\end{theorem}
\begin{proof}
For each admissible state, the prior order comparison gives
\(\N_\lambda(\sigma)\succeq c\sigma\), where
\(c=\lambda+(1-\lambda)/\kappa>0\).
Indeed \(1+\lambda(M_d-1)\ge d/[2(d^2-1)]>0\).
An equal-reference difference obeys \(\N_\lambda(D)=\lambda D\). Inverse antitonicity for \((L_\sigma+R_\sigma)/2\), or for \(L_\rho+uR_\sigma\) using both state budgets, gives the factor \(\lambda^2/c\). Include the quadratic endpoint and the positive integral in \eqref{eq:resolvent}. The new lower, rather than this comparison machinery, is the point of the construction.

Put \(a=d-1\), \(r=2a+1\), \(q=a+e^2\), \(e>0\), and use the reference labels \(0,\ldots,2a\). Define
\[
 v_0=\sum_{j=0}^{a-1}|j+1,j\rangle+e|0,2a\rangle,\qquad
 v_1=\sum_{j=0}^{a-1}|j,a+j\rangle+e|d-1,2a\rangle.
\]
Let \(b_0=|0,2a\rangle\), \(b_1=|d-1,2a\rangle\), and
\[
 \sigma_e=\frac{v_0v_0^*+v_1v_1^*}{2q},\qquad
 D=v_0b_0^*+b_0v_0^*-v_1b_1^*-b_1v_1^*.
\]
The supports are disjoint, \(\|v_i\|^2=q\), and \(Q=2\sigma_e\) is a projection. All principal--center cross terms change the A index. The two diagonal center derivatives cancel in the same reference label. Thus \(\Tr_AD=0\).

The actual reference marginal has mass \(1/(2q)\) on its \(2a\) principal labels and \(e^2/q\) on the center. Consequently \(B_e=\E_A(\sigma_e)>0\) and
\[
 \frac{[v_0,v_1]^*B_e^{-1}[v_0,v_1]}{2q}
 =d(a+1/2)I_2=M_d I_2.
\]
This proves that the whitened state has two nonzero eigenvalues \(M_d\). Also \((I-Q)D(I-Q)=0\); its boundary SLD form, evaluated spectrally without inverting the singular base, is \(16q\).

Set \(s_{e,p}=(1-p)B_e+p\sigma_e\). Each support has an invariant plane spanned by its uniform principal vector and center coordinate. In orthonormal coordinates the two base blocks are
\begin{equation}\label{eq:principal-plane}
 W_p=\frac1{2dq}
 \begin{pmatrix}A_p&dpe\sqrt a\\dpe\sqrt a&C_pe^2\end{pmatrix},
 \quad T=\begin{pmatrix}0&\sqrt a\\\sqrt a&2e\end{pmatrix},
\end{equation}
where \(A_p=1+(da-1)p\), \(C_p=2+(d-2)p\); the tangent blocks are \(T,-T\).
Each support has \(a-1\) additional principal modes with eigenvalue \((1-p)/(2dq)\) and zero tangent. Outside these supports the remaining \(d(2d-1)-2d\) coordinates are diagonal replacement modes with zero tangent. This accounts for the full dimension.
The aligned determinant is
\[
 \det W_p=\frac{e^2(1-p)[1+(M_d-1)p]}{2d^2q^2}>0
 \qquad\left(-\frac1{M_d-1}<p<1\right).
\]
Both \(p=\theta\) and \(p=\lambda\theta\) lie in this domain for \(0<\theta<1\) and the entire physical range of \(\lambda\).

For \(W_pY+YW_p=2T\), the two-dimensional identity
\[
 Y=\frac{T+\det(W_p)W_p^{-1}TW_p^{-1}}{\Tr W_p}
\]
follows by Cayley--Hamilton and direct substitution. Inserting \eqref{eq:principal-plane} and taking the limit at fixed admissible \(p\) gives
\begin{equation}\label{eq:general-limit}
 \lim_{e\downarrow0}\Tr(TY)=\frac{8aM_d}{1+(M_d-1)p}.
\end{equation}
The full ambient SLD form is twice this quantity. At fixed \(\theta<1\), its retention therefore tends to
\(\lambda^2[1+(M_d-1)\theta]/[1+\lambda(M_d-1)\theta]\).

Choose \(0<\theta<(\kappa-1)/(M_d-1)\). Both \(s_{e,\theta}\) and \(\kappa B_e-s_{e,\theta}\) are positive definite, so sufficiently small bilateral perturbations \(\rho_t=s_{e,\theta}+tD\) preserve faithfulness, the actual marginal and both budgets. For triangular divergence take \(t\to0\) first; for minimal chi-square the ratio is exactly quadratic. Then take \(e\downarrow0\) at fixed \(\theta\), finally \(\theta\uparrow(\kappa-1)/(M_d-1)\). This reaches \eqref{eq:profile} as a supremum. The endpoints \(\lambda=0,1\) are immediate. No exchange of singular limits or finite faithful maximizer is needed.
\end{proof}

The earlier rank-two allocation achieved the smaller cap \(L_d=d(d+1)/2\) with reference dimension \(d+1\). The principal-rich allocation raises the cap by \(d(d-2)/2\), at the cost of reference dimension \(2d-1\); for \(d=2\) they coincide. This is a development of the same mechanism, not an independently invented theory or a free reference-size improvement.

\begin{remark}[The empty endpoint and the domain of the order bound]
At \(\kappa=1\), positivity and equal traces force \(\rho=\E_A(\rho)\) and \(\sigma=\E_A(\sigma)\); equal reference then forces equality. There are no distinct finite pairs, so \eqref{eq:profile} is not assigned to that empty optimization. The base-only SLD problem with unrestricted tangents at a product base is a different problem. Likewise the expression \(\lambda^2/c\) must never be used when \(c\le0\); for example \(\lambda=-1/3,\kappa=5\) gives a negative raw expression in dimension two, not a contraction bound.
\end{remark}

\section{Exact qubit resource restrictions}\label{sec:resources}
For \(d=2\), Theorems~\ref{thm:pauli} and~\ref{thm:order} combine to give, for every \(\kappa>1\),
\[
 C_\kappa(\lambda)=\frac{\lambda^2K}{1+\lambda(K-1)},\qquad K=\min(\kappa,3),\qquad -1/3\le\lambda\le1.
\]
For \(\kappa\ge3\), the upper is the unrestricted isotropic case of Theorem~\ref{thm:pauli}, \(3\lambda^2/(1+2\lambda)\); the same lower family has order arbitrarily close to three. At the channel endpoints use DPI or the identity directly.

For a qubit and any finite reference,
\[
 2\E_A(\sigma)-\sigma=(\sigma_y\otimes I)\sigma^{T_A}(\sigma_y\otimes I).
\]
Thus PPT bases obey the order-two bound, as do separable bases. The exact complete SLD/minimal-chi-square/triangular coefficient restricted to PPT or separable models is
\(2\lambda^2/(1+\lambda)\). For finite Petz both states must belong to the class.
To prove the matching lower, use the qubit ladder above (equivalent under a local basis swap to the earlier allocation) with \(s_{e,\theta}\) and \(\theta<1/2\). Its partial transpose has two blocks with determinant
\(e^2(1-\theta-2\theta^2)/[8(1+e^2)^2]>0\) and two positive singletons. Small bilateral perturbations stay PPT. In dimension \(2\times3\), the prior Horodecki theorem~\cite{HHH} converts these pairs to separable ones. This uses PPT=SEP only for the lower in \(2\times3\), not for arbitrary references or larger local dimensions.

For \emph{each} normalized nonaffine operator-convex Petz generator, the SEP/PPT pair coefficient is also exactly \(2\lambda^2/(1+\lambda)\). Indeed, the order-two comparison gives the upper, while the qubit Bell family of Proposition~\ref{prop:bell}, with \(K<2\), has all Bell probabilities \(p_j\) strictly below \(1/2\). The partial-transpose eigenvalues are \(1/2-p_j\). Both the base and sufficiently small finite perturbations are strictly PPT, hence separable in \(2\times2\), and approach the stated lower. The same gives the budget profile inside SEP/PPT for \(1<\kappa\le2\). This restricted-pair statement includes Umegaki but does not determine the unrestricted coefficient or actual MI, whose second state has an additional product-marginal requirement.

A classical-reference \emph{family} has both states and the tangent block diagonal in one fixed reference basis, with fixed marginal weights. The forms split into weighted ordinary-qubit forms, whose coefficient is \(\lambda^2\) by~\cite{HR}; \(R=1\) gives equality. Classicality of the base alone is insufficient. For example at \(\lambda=-1/3\), the classical base \(\sigma=\operatorname{diag}(9,1,1,9)/20\) with \(D_{03}=D_{30}=1\) gives retention \(3/11>1/9\); the direction creates reference coherences.

At \(\lambda=-1/3\) the exact levels are therefore
\[
 \text{classical-reference family}:\frac19,\qquad
 \text{SEP/PPT family}:\frac13,\qquad
 \text{unrestricted quantum family}:1.
\]
SLD retention (or minimal-chi-square retention with base \(\sigma\)) above \(2\lambda^2/(1+\lambda)\) is a sufficient certificate that the input base is NPT. For a general finite Petz pair, the proved implication is only that exceeding this threshold excludes both input states being PPT; it does not identify which state is NPT. This mathematical implication is a contrapositive of the prior reduction criterion plus the quantitative upper; it is not an experimentally demonstrated protocol.
